\centerline{\bf \Huge Домашнее задание}
\centerline{\bf \Huge Теория игр}

\begin{flushright}
        \scriptsize{У нас нет времени на игры --- мы должны воевать!}
\end{flushright}

\problem{1!} Внутрь прямоугольника размерами $10\times100$ Герман и Алина по очереди кладут палочки длины 10. Палочки должны находиться строго внутри прямоугольника и не пересекаться с другими. Считается, что у палочек есть какая-то маленькая толщина. Проигрывает тот, кто не сможет сделать ход. Алина начинает первая.

\problem{2!} Герман и Дольган попали в гендзюцу Ердема под названием шахматное цукуёми. В ряд выложены 2002 пешки, пронумерованные от 1 до 2002. За один ход каждый из них может менять местами две пешки находящиеся через одну друг от друга. Смогут ли они договориться так, чтобы спустя несколько ходов расположить пешки в обратном порядке или они навсегда застряли в гендзюцу Ердема?

\problem{3} В кастрюле лежит 2027 беригов. Герман и Ердем по очереди съедают из кастрюли от 1 до 17 беригов. Когда кастрюля опустошается, игроки подсчитывают съеденное количество беригов. Если эти числа взаимно просты – выигрывает Ердем, в противном случае выигрывает Герман. Ердем ходит первым.

\problem{4} Герман и Анджур Аркадьевич играют в следующую игру. Герман выписывает подряд цифры 0 или 1 в произвольном порядке, пока не получится 2023-значное число. Каждый раз, когда Герман выставляет цифры, Анджур Аркадьевич меняет местами две уже написанные цифры(первый ход он пропускает, а остальные обязательно ходит!) Если итоговое число получится палиндромом(то есть будет читаться с обеих сторон одинаково), то Анджур Аркадьевич побеждает, если нет, то побеждает Герман.

\problem{5} Конь стоит на поле $a1$. За ход разрешается передвигать коня на две клетки вправо и одну клетку вверх или вниз, или на две вверх и на одну вправо или влево. Проигрывает тот, кто не может сделать ход.